% Generated by scripts/assemble.py; corrections belong in source/corrections.json.
\examhead{78}{东北师范大学 2024 年数学分析试卷}{76}
% p077-b22; source PDF page 77
\par\Needspace{5\baselineskip}
\blockanchor{p077-b22}
\majorquestion{一}计算题. 每题6分, 共30分.\par
% p077-b23; source PDF page 77
\blockanchor{p077-b23}
(1)  $\lim_{n\to\infty}\left(\frac{1}{n}+\frac{1}{n+1}+\cdots+\frac{1}{2n+1}\right).$\par
% p077-b24; source PDF page 77
\blockanchor{p077-b24}
(2) $\lim_{x\to +\infty}\sqrt{x}\int_{x}^{x + 1}\frac{\mathrm{d}t}{\sqrt{t + \ln t}}.$\par
% p077-b25; source PDF page 77
\blockanchor{p077-b25}
(3)  $\lim_{n\to\infty}\left(\frac{1}{2}+\frac{3}{2^{2}}+\frac{5}{2^{3}}+\cdots+\frac{2n-1}{2^{n}}\right).$\par
% p078-b00; source PDF page 78
\blockanchor{p078-b00}
(4) $\int_0^1\mathrm{d}x\int_0^1\left|xy - \frac{1}{4}\right|\mathrm{d}y.$\par
% p078-b01; source PDF page 78
\par\Needspace{6\baselineskip}
\blockanchor{p078-b01}
(5) 求球面 $x = \sin \varphi \cos \theta, y = \sin \varphi \sin \theta, z = \cos \varphi$ 在对应 $\theta = \varphi = \frac{\pi}{4}$ 处的切平面和法线方程.
\majorquestion{二}（15分）由方程组\par
\nopagebreak[4]
% p078-b02; source PDF page 78
\blockanchor{p078-b02}
\displayformula{\left\{ \begin{array}{l} x ^ {2} + y ^ {2} - u v = 0; \\ x y + u ^ {2} - v ^ {2} = 0. \end{array} \right.}
% p078-b03; source PDF page 78
\blockanchor{p078-b03}
能否确定 u, v 为 x 与 y 的函数？在能确定隐函数的条件下，求  $u_{x}, v_{x}, u_{y}, v_{y}$ .\par
% p078-b04; source PDF page 78
\par\Needspace{5\baselineskip}
\par\Needspace{6\baselineskip}
\blockanchor{p078-b04}
\majorquestion{三}(15 分) 求函数 $f(x, y) = \sin x + \sin y + \sin (x + y)$ 在 $(x, y) \in [0, 2\pi] \times [0, 2\pi]$ 的所有极值点.
\majorquestion{四}(15 分) 计算曲面积分\par
\nopagebreak[4]
% p078-b05; source PDF page 78
\blockanchor{p078-b05}
\displayformula{I = \iint_ {S} (z ^ {2} + x) \mathrm{d} y \mathrm{d} z + \sqrt {z} \mathrm{d} x \mathrm{d} y.}
% p078-b06; source PDF page 78
\blockanchor{p078-b06}
其中 $S$ 为抛物面 $z = \frac{1}{2} (x^2 + y^2)$ 在 $z = 0$ 和 $z = \pi$ 之间的部分, 定向取下侧.\par
% p078-b07; source PDF page 78
\par\Needspace{5\baselineskip}
\par\Needspace{6\baselineskip}
\blockanchor{p078-b07}
\majorquestion{五}（15分）设函数 $f(x)$ 在 $(- \infty, + \infty)$ 上有任意阶导数，且对任意实数 $x$ 及 $n = 0,1,2,\dots$ ，满足\par
\nopagebreak[4]
% p078-b08; source PDF page 78
\blockanchor{p078-b08}
\displayformula{| f ^ {(n)} (x) | \leq n! | x |.}
% p078-b09; source PDF page 78
\blockanchor{p078-b09}
证明： $f(x) = 0$ ， $x\in (-\infty , + \infty)$\par
% p078-b10; source PDF page 78
\par\Needspace{5\baselineskip}
\blockanchor{p078-b10}
\majorquestion{六}设 $D$ 为平面区域, $u(x, y) \in C^2(D)$, 证明: $u(x, y)$ 为调和函数, 即 $u$ 满足 $\frac{\partial^2 u}{\partial x^2} + \frac{\partial^2 u}{\partial y^2} = 0$ 的充要条件是: 对 $D$ 内任一圆周 $L$, 且 $L$ 所围圆属于 $D$, 都有 $\int_{L} \frac{\partial u}{\partial n} \mathrm{d}s = 0$.\par
% p078-b11; source PDF page 78
\par\Needspace{5\baselineskip}
\blockanchor{p078-b11}
\majorquestion{七}(15 分) 设 $\{a_{n}\}$ 是正的单调递增序列, 证明: 级数 $\sum_{n=1}^{\infty} \left(1 - \frac{a_{n}}{a_{n+1}}\right)$ 当 $\{a_{n}\}$ 有界时收敛, 当 $\{a_{n}\}$ 无界时发散.\par
% p078-b12; source PDF page 78
\par\Needspace{5\baselineskip}
\blockanchor{p078-b12}
\majorquestion{八}(15 分) 令 $f(t) = \int_{1}^{+\infty} \frac{\cos xt}{1 + x^2} \, \mathrm{d}x$. 证明:\par
% p078-b13; source PDF page 78
\blockanchor{p078-b13}
(1) 积分在 $t \in (-\infty, +\infty)$ 上一致连续. (按道理应该是一致收敛)\par
% p078-b14; source PDF page 78
\blockanchor{p078-b14}
(2) $f(t)\in C(-\infty , + \infty).$\par
% p078-b15; source PDF page 78
\blockanchor{p078-b15}
(3)  $\lim_{t\to+\infty}f(t)=0.$\par
% p078-b16; source PDF page 78
\par\Needspace{5\baselineskip}
\blockanchor{p078-b16}
\majorquestion{九}(15 分) 设 $\{f_n\}$ 是 $[a, b]$ 上连续的函数列, 对任意的 $x \in [a, b]$, 有 $f_n(x) \leq f_{n+1}(x)$, $n = 1, 2, \cdots$. 证明: 如果 $\{f_n\}$ 在 $[a, b]$ 上收敛于连续函数 $f$, 那么 $\{f_n\}$ 必定一致收敛于 $f$.\par
